\begin{enumerate}
    \item On dispose de barreaux aimantés identiques créant chacun en un point
          $M$, équidistant des deux aimants, des champs  magnétiques notés
          respectivement $\vec{B}_{1}$ et $\vec{B}_{2}$.

          \begin{figure}[H]
              \centering
              \begin{tikzpicture}
                  \def\d{2}
                  \foreach \x/\y/\a [count=\i] in {-\d/0/east,\d/0/west}{%
                  \node[inner sep=0pt,draw,thick,minimum width=2cm,
                      minimum height=6mm,
                      anchor=\a,label=(\i)] (A\i) at (\x,\y) {};
                  \draw (A\i.north) -- (A\i.south);
                  }
                  \node[circle,inner sep=1pt,fill,
                      label={[above=1mm]:$M$}] (M) at (0,0) {};
                  \draw[densely dashed] (A1.east) -- (M) -- (A2.west);
                  \node[left] at (A1.east) {S};
                  \node[right] at (A1.west) {N};
                  \node[right] at (A2.west) {S};
                  \node[left] at (A2.east) {N};
                  \begin{scope}[xshift=8cm,yshift=0cm]
                      \def\al{35}
                      \foreach \y/\a [count=\i] in
                          {1.3/180-0.5*\al,-1.3/180+0.5*\al}{%
                      \node[inner sep=0pt,draw,thick,minimum width=2cm,
                          minimum height=6mm,
                          rotate=\a,label={[below,rotate=180+\a]:(\i)}]
                          (A\i) at (-2,\y) {};
                      \draw (A\i.north) -- (A\i.south);
                      }
                      \path[name path=D1] (A1.west) -- ++(-0.5*\al:5);
                      \path[name path=D2] (A2.west) -- ++(0.5*\al:5);
                      \fill[name intersections={of=D1 and D2,by=M}]
                          (M) circle (1.4pt)
                          node[above=1mm] {$M$};
                      \draw[densely dashed] (A1.west) -- (M) -- (A2.west);
                      \node[rotate=-0.5*\al,anchor=180] at (A1.east) {N};
                      \node[rotate=-0.5*\al,anchor=0] at (A1.west) {S};
                      \node[rotate=0.5*\al,anchor=180] at (A2.east) {N};
                      \node[rotate=0.5*\al,anchor=0] at (A2.west) {S};
                  \end{scope}
              \end{tikzpicture}
          \end{figure}
          \begin{enumerate}
              \item Représenter le vecteur champ magnétique crée par chaque
                    aimant au point $M$, dans chaque cas.~\textbf{(1pt)}
              \item Représenter le vecteur champ magnétique résultant au point
                    $M$, dans chaque cas. \textbf{(1pt)}
          \end{enumerate}
    \item Le spectre magnétique d’un aimant droit est représenté sur le schéma
          ci-dessous.

          \begin{minipage}[t][][t]{0.44\linewidth}
          \begin{enumerate}
              \item Orienter les lignes de champ représentées. \textbf{(0,5pt)}
              \item Tracer le vecteur champ magnétique aux points $M_{i}$.
                    \textbf{(1pt)}
              \item Dessiner une aiguille aimantée en chaque point $M_{i}$, en
                    précisant la nature de ses pôles. \textbf{(1pt)}
          \end{enumerate}
          \end{minipage}
          \hfill
          \begin{minipage}[t][][t]{0.54\linewidth}
          \vspace{-5mm}
          \begin{figure}[H]
              \centering
              \begin{tikzpicture}
                  \clip (0,0) ellipse (4.5cm and 2.5cm);
                  \node[inner sep=0pt,draw,thick,minimum width=5cm,
                      minimum height=1cm] (A) {};
                  \draw (A.north) -- (A.south);
                  \foreach \x [count=\i] in {0.05,0.2,0.35}{%
                      \draw ($(A.north west)!\x!(A.north east)$) to[bend left=70]
                          node[pos=0.5] (M\i) {} 
                          ($(A.north west)!1-\x!(A.north east)$);
                      \draw ($(A.south west)!\x!(A.south east)$) to[bend right=70]
                          node[pos=0.75] (N\i) {} 
                          ($(A.south west)!1-\x!(A.south east)$);
                  }
                  \foreach \x/\a/\l [count=\i] in {0.05/130/2.5,0.3/150/4}{%
                      \draw ($(A.north west)!\x!(A.south west)$)
                          to[bend left=\a,looseness=\l]
                          node[pos=0.9] (P\i) {} 
                          ($(A.north east)!\x!(A.south east)$);
                      \draw ($(A.north west)!1-\x!(A.south west)$)
                          to[bend right=\a,looseness=\l]
                          ($(A.north east)!1-\x!(A.south east)$);
                  }
                  \draw (A.east) -- ++(2,0);
                  \draw (A.west) -- ++(-2,0)
                      node[pos=0.6,circle,inner sep=1pt,fill,label=$M_{2}$] {};
                  \node[left=2mm,font=\huge] at (A.east) {N};
                  \node[right=2mm,font=\huge] at (A.west) {S};
                  % Mi
                  \node[circle,inner sep=1pt,fill,label=$M_{1}$] at (M1) {};
                  \node[circle,inner sep=1pt,fill,label=$M_{3}$] at (N2) {};
                  \node[circle,inner sep=1pt,fill,label=left:$M_{4}$] at (P1) {};
              \end{tikzpicture}
          \end{figure}
          \end{minipage}

    \item 
          \begin{minipage}[t][][t]{0.50\linewidth}
    Représenter le vecteur champ magnétique $\vec{B}$ créé par le conducteur aux
    points $M_{1}$ et $M_{2}$.~\textbf{(0,5)}

    En déduire le sens du courant $I$.~\textbf{(0,5)}
\end{minipage}
\hfill
\begin{minipage}[t][][t]{0.48\linewidth}
    \begin{figure}[H]
        \centering
        \begin{tikzpicture}
            \def\d{5}
            \def\ar{1.6}
            % Aiguille
            \node[diamond,inner sep=0pt,minimum width=7mm,
                  minimum height=3.4mm] (aa) at (279:\ar) {};
            \path[draw,fill=white,rounded corners=0.4pt]
                (aa.north) -- (aa.south) -- (aa.east) -- cycle;
            \path[draw,fill=red,rounded corners=0.4pt]
                (aa.north) -- (aa.south) -- (aa.west) -- cycle;
            \node[below,inner sep=4pt,font=\scriptsize] at (aa.east) {S};
            \node[below,inner sep=4pt,font=\scriptsize] at (aa.west) {N};
            \draw[] (0,0) circle (\ar);
            \draw[<-,shorten <=2pt] (180:\ar) -- ++(-1,0)
                node[left,align=center,font=\scriptsize] {Ligne de\\ champ};
            % M
            \node[circle,fill,inner sep=1pt,
                label={[above]:$M_{1}$}] (M) at (aa.center) {};
            \node[circle,fill,inner sep=1pt,
                label={[above right=-1pt,inner sep=0pt]:$M_{2}$}] at (45:\ar) {};
            % fil
            \node[circle,fill=gray,inner sep=1pt,
                label={[above,inner sep=4pt,font=\scriptsize,
                align=center]:{fill\\ rectiligne}}] at (0,0) {};
        \end{tikzpicture}
    \end{figure}
\end{minipage}


	



\end{enumerate}

